\section{Infinite-volume spectral type}
\label{sec:spectral}

We now remove the fixed-energy restriction from the spectral conclusion.
The probabilistic input is the true finite-torus estimate
\cref{boot:fractional}. We use the bounded self-adjoint spectral theorem
and its Borel functional calculus in the form of
\cite[Theorems~3.1 and~3.7]{Teschl2009}. We include the spectral-averaging
step to make the quantifiers explicit; it is the rank-one mechanism of
\cite[Theorem~2', Proposition~7 and Theorem~8(i)]{SimonWolff1986}.

\begin{lemma}\label{spec:column}
For every fixed $h>0$, every fixed $E\in(-4-h,4+h)$ and every site
$x\in\Z^2$, almost surely there is $u\in\ell^2(\Z^2)$ such that
\[
 (H_v-E)u=\delta_x.
\]
If $\mu_x$ is the scalar spectral measure of $\delta_x$, then almost
surely
\begin{equation}\label{spec:gamma}
 \Gamma_x(E):=\int_\R\frac{\mu_x(\dd e)}{(e-E)^2}<\infty.
\end{equation}
The integrand is interpreted as $+\infty$ at $e=E$.
\end{lemma}

\begin{proof}
Couple the torus of side $b_N$ to the infinite sample by taking the true
potentials in the square of that side centered at $x$. Let $G_N(E)$ be
its real Green matrix. By \cref{boot:fractional}, there are fixed
$q,c,C>0$, with $q<1$, such that for all sufficiently large $N$,
\[
 \E\abs{G_N(E)_{yx}}^q\le C\exp(-c\abs{y-x}_\infty).
\]
The torus distance from $x$ equals the physical coordinate distance in
this centered square. Hence for $0<a<c$,
\[
 \sup_N\E\sum_{y\text{ in the square}}
 e^{a\abs{y-x}_\infty}\abs{G_N(E)_{yx}}^q<\infty,
\]
where the supremum may omit the finitely many smaller tori.
Fatou implies that the liminf of these nonnegative sums is finite
almost surely. Choose a subsequence on which they are bounded. Its
coordinates are bounded, and a diagonal subsequence converges at every
lattice site. Fatou for the counting measure gives a limit $u$ with
finite weighted $\ell^q$ sum. In particular $u\in\ell^2$: an $\ell^q$
sequence is bounded, and
$\sum_y\abs{u_y}^2\le\norm{u}_\infty^{2-q}\sum_y\abs{u_y}^q$.
Every fixed nearest-neighbor equation eventually avoids the torus
boundary, so the limit solves $(H_v-E)u=\delta_x$.

The bounded self-adjoint spectral theorem gives
\[
 \mu_x(\dd e)=(e-E)^2\mu_u(\dd e),
\]
where $\mu_u$ is the scalar spectral measure of $u$. Thus
$\mu_x(\{E\})=0$ and the integral in \cref{spec:gamma} is at most
$\norm{u}^2$.
\end{proof}

\begin{lemma}[Rank-one spectral averaging]\label{spec:averaging}
Let $H_0$ be bounded self-adjoint and let $\xi$ be a unit vector.
For $a\in\R$, let $\mu_a$ be the scalar spectral measure of $\xi$
for $H_0+a\langle\xi,\cdot\rangle\xi$. Then, as measures on $\R$,
\begin{equation}\label{spec:average}
 \int_\R\mu_a\dd a\le \operatorname{Leb}.
\end{equation}
Suppose on a Borel set $B$ the base measure $\mu_0$ satisfies
$\int(e-E)^{-2}\mu_0(\dd e)<\infty$ for Lebesgue-almost every
$E\in B$. For almost every $a$, the restriction of $\mu_a$ to $B$
is concentrated on eigenvalues of the perturbed operator.
\end{lemma}

\begin{proof}
For $\operatorname{Im}z>0$ write
\[
 F(z)=\int\frac{\mu_0(\dd e)}{e-z},\qquad
 F_a(z)=\frac{F(z)}{1+aF(z)}.
\]
The second identity follows from rank-one resolvent inversion. Since
$\operatorname{Im}F(z)>0$, direct integration in $a$ gives
$\int_\R\operatorname{Im}F_a(z)\dd a=\pi$.
For a nonnegative continuous compactly supported function $f$, integrate
against $f(\operatorname{Re}z)$, use Tonelli, and let
$\operatorname{Im}z\downarrow0$. The normalized Poisson kernels are an
approximate identity, and Fatou in $a$ gives
\[
 \int_\R\int_\R f(e)\mu_a(\dd e)\dd a
 \le\int_\R f(e)\dd e.
\]
This proves \cref{spec:average}.

Let $B_0$ be the subset of $B$ where the base inverse-square integral
is finite. It is Borel, and $B\setminus B_0$ is Lebesgue null.
Equation~\eqref{spec:average} shows that it has zero $\mu_a$-measure
for almost every $a$. At $E\in B_0$, Cauchy--Schwarz and dominated
convergence show that $F(E+i0)$ exists and is real. If
$1+aF(E+i0)\ne0$, then
\[
 \operatorname{Im}F_a(E+i\eta)\longrightarrow0,
 \qquad
 \frac{\mu_a([E-\eta,E+\eta])}{\eta}
 \le2\operatorname{Im}F_a(E+i\eta)\longrightarrow0.
\]
For any finite Borel measure $\mu$, the set of points where this last
ratio tends to zero has $\mu$-measure zero. To see this directly on a
bounded interval, fix $\varepsilon>0$ and restrict first to points where
the ratios are at most $\varepsilon$ at every dyadic radius below a
common threshold. At a finer dyadic mesh, each interval meeting this
set lies in a centered interval of that radius at one of its points,
so its measure is at most $\varepsilon$ times the mesh length. Summing
gives at most a constant times $\varepsilon$ times the length of the
bounded ambient interval. Increase the threshold sets to their union,
then let $\varepsilon\downarrow0$, and finally exhaust the real line.

It remains to consider $E\in B_0$ with $1+aF(E+i0)=0$. For each
positive integer $k$, apply the bounded Borel multiplier
$\one_{\{|e-E|\ge1/k\}}(e-E)^{-1}$ for $H_0$ to $\xi$.
The inverse-square integral makes these vectors Cauchy in norm.
Their limit $u$ satisfies
\[
 (H_0-E)u=\xi,\qquad \langle\xi,u\rangle=F(E+i0).
\]
Indeed, the base spectral measure has no atom at $E$, so the truncated
equations converge to $(H_0-E)u=\xi$ by boundedness of $H_0$.
Cauchy--Schwarz gives the asserted inner product. Consequently
\[
 (H_0+a\langle\xi,\cdot\rangle\xi-E)u
       =(1+aF(E+i0))\xi=0.
\]
The vector is nonzero because its base equation has a nonzero right
side. Hence $E$ is an eigenvalue. All these eigenvectors belong to the
closed cyclic subspace generated by $\xi$ under $H_0$. This subspace is
separable and invariant under the perturbed operator. Eigenvectors for
distinct eigenvalues are orthogonal, so the energies obtained in this
way form a countable set. The remaining scalar spectral mass is
therefore pure point, as asserted.
\end{proof}

\begin{proposition}\label{spec:readout}
For each fixed $h>0$, almost surely every site spectral measure of
$H_v$ is pure point. Consequently $H_v$ has pure-point spectral type
and a complete orthonormal eigenbasis.
\end{proposition}

\begin{proof}
The property in \cref{spec:gamma} is jointly measurable in the true
potential configuration and $E$. Indeed $\norm{H_v}\le4+h$, and
polynomial functional calculus followed by uniform approximation of
continuous functions makes $v\mapsto\mu_x$ a Borel measure kernel.
Integration of the nonnegative Borel kernel $(e-E)^{-2}$ then gives
the asserted joint measurability. This property no longer mentions
any auxiliary port construction or its energy-dependent parameters.

By \cref{spec:column} and Fubini, for almost every setting of the
potentials away from $x$, almost every base value $v_x^0\in[-h,h]$
has finite $\Gamma_x(E)$ for Lebesgue-almost every
$E\in(-4-h,4+h)$. Fix such an exterior and one such base value.
Varying $v_x$ corresponds to the rank-one parameter $a=v_x-v_x^0$.
By \cref{spec:averaging}, for almost every allowed $a$ the site
measure is pure point on that open interval. The same averaging
inequality removes the two endpoint energies. For every allowed $a$
the norm bound confines the spectrum to the closed interval, so the
entire site measure is pure point.

This conditional conclusion integrates to an almost-sure assertion.
For completeness, purity of a scalar measure kernel is measurable:
$(v,e)\mapsto\mu_x(v)(\{e\})$ is Borel by continuous approximations
to the diagonal indicator, and its zero set has measurable
$\mu_x(v)$-mass. Purity is exactly the vanishing of that mass.
Intersect now over the countable set of sites $x$. Each $\delta_x$
lies in the closed pure-point subspace; their span is dense, so that
subspace is the entire Hilbert space. Choosing orthonormal bases in
the mutually orthogonal eigenspaces yields the required eigenbasis.
We have not assumed that one site vector is cyclic for the full
two-dimensional operator.
\end{proof}

\begin{proposition}\label{spec:spectrum}
For each fixed $h>0$, almost surely
$\sigma(H_v)=[-4-h,4+h]$. Under \cref{spec:readout}, the eigenvalues
are dense in this interval.
\end{proposition}

\begin{proof}
The norm bound gives the inclusion into the displayed interval.
For each prescribed constant in a countable dense subset of $[-h,h]$,
each positive rational accuracy, and each integer box size, infinitely
many independent disjoint-box trials have positive probability that
every potential in the box is within that accuracy of the constant.
Almost surely such boxes therefore occur; take the countable
intersection over these choices.

The free adjacency has plane waves with energies
$2\cos k_1+2\cos k_2$, whose range is $[-4,4]$. Normalize a plane
wave cut off to one of the large nearly constant boxes. The hopping
error is supported within bounded distance of its boundary and has
norm tending to zero because boundary size divided by box volume
tends to zero. The diagonal error is at most the prescribed accuracy.
Let size tend to infinity and accuracy to zero, using countable dense
choices of the wave energies and the constants. The approximate
eigenvector criterion puts a dense subset of $[-4-h,4+h]$ in the
spectrum, and its closedness gives the entire interval.

For an operator with a complete orthonormal eigenbasis, the spectrum
is the closure of its eigenvalues: outside that closure the reciprocal
of the diagonal eigenvalue differences defines a bounded inverse.
The last assertion follows. Together with \cref{spec:readout}, this
proves \cref{main:theorem}.
\end{proof}
